\scanpage{5}
\uline{Arbres plans.}\hfill\mspage{1}

Rappelons qu'un \uline{arbre}\inserted{(comb.)} est un graphe (comb.) 1-connexe (ce qui implique que c'est un graphe strict, i.e. sans \unsure{bords|boucles} fermés, et pas deux \deleted{\unsure{pts}} sommets distincts joints \unsure{au|en} plus une arête). On peut donc considérer la structure donnée par $(S,A)$, $S$ ens des sommets, $A\subset\mathcal P_2(S)$. \deleted{$\unsure{A}$} L'ens $S$ est $\ne\varnothing$, et $A\deleted{\ne}=\varnothing\iff\operatorname{card}S=1$ (arbre \uline{ponctuel}).

On appelle \uline{arbre plan} un arbre muni d'une classe de plongements\inserted{(de sa réal. top.)} dans une surface (deux plongements étant identifiés, $|\Gamma|\hookrightarrow X$, $|\Gamma|\hookrightarrow X'$ s'il existe des voisinages\inserted{$U,U'$} de $|\Gamma|$ ds $X,X'$ et un homéo. $U\simeq U'$ induisant l'identité sur $|\Gamma|$). Les \inserted{classes de} plongements (vu que $|\Gamma|$ est 1-connexe et admet un voisinage ouvert 1-connexe $U$ dans $X$; -- d'où $U\simeq\unsure{\mathbb R|\mathbb R^2}$ --) s'identifient aux \uline{bicircularisations} de $\Gamma$, i.e. les \deleted{applications} paires $(\varepsilon,\unsure{c|\varphi})$, $\varepsilon$ un ens à 2 éléments, \deleted{$\unsure{c}$} et $c:\varepsilon\longrightarrow\operatorname{circ}(\Gamma)$, $c$ compatible aux ``passages à l'opposé''. Ici $\operatorname{circ}(\Gamma)$ est l'ens des \uline{circularisations} de $\Gamma$, i.e. des systèmes d'ordres circulaires sur les $A_s$ (où pour $\forall s\in S$, $A_s$ est l'ens des arêtes incidentes à $s$). Une \inserted{classe de} plongement de $|\Gamma|$ dans une surface \uline{orientée} correspond à une circularisation de $\Gamma$; ça ne dépend que de la famille des $A_s$ $(s\in S)$ tels que $\operatorname{card}A_s\geq3$. Donc si les sommets de $\Gamma$ sont d'ordre $\leq\unsure{1|2}$, il y en a \unsure{un} et un seul
\mspage{4}
